\documentclass[a4paper,12pt]{article}
\usepackage[utf8]{inputenc}
\usepackage[T1]{fontenc}
\usepackage[french]{babel}
\usepackage{geometry}
\usepackage{amsmath}
\usepackage{amssymb}
\usepackage{enumitem}
\usepackage{booktabs}
\usepackage{xcolor}
\usepackage{fancyhdr}
\usepackage{tcolorbox}
\usepackage{hyperref}

\geometry{margin=2.5cm}
\pagestyle{fancy}
\fancyhead[L]{\textsc{ImmoGest v1.0}}
\fancyhead[R]{\textsc{Logique de Calcul}}
\fancyfoot[C]{\thepage}

\definecolor{formule}{RGB}{24,78,119}
\definecolor{exemple}{RGB}{46,125,50}

\newtcolorbox{formulebox}[1][]{colback=formule!5, colframe=formule!80, title=#1, fonttitle=\bfseries}
\newtcolorbox{exemplebox}[1][]{colback=exemple!5, colframe=exemple!70, title=#1, fonttitle=\bfseries}

\title{\vspace{-1cm}\textbf{\LARGE ImmoGest — Logique de Calcul} \\ \large Référentiel des formules de l'application}
\author{IMMOSTAR SCI}
\date{Juin 2026}

\begin{document}
\maketitle
\tableofcontents
\newpage

% =================================================================
\section{Montant Attendu (Multi-Baux)}
% =================================================================

Le montant attendu représente la totalité des loyers et charges qu'un locataire aurait dû payer depuis son entrée dans le logement. Pour un locataire ayant eu \textbf{plusieurs baux successifs} sur le même appartement, chaque bail contribue séparément.

\begin{formulebox}[Formule générale]
\begin{equation}
\text{Attendu total} = \sum_{\text{chaque bail}} \text{Attendu}_{\text{bail}}
\end{equation}
\end{formulebox}

\subsection{Bail MENSUEL}

\begin{formulebox}[Attendu — Bail mensuel]
\begin{equation}
\text{Attendu} = (\text{loyer} + \text{charges}) \times \frac{\text{jours}}{30}
\end{equation}
\vspace{-0.3cm}
\begin{itemize}[noitemsep]
  \item Bail terminé : jours = dateFin $-$ dateDebut
  \item Bail actif : jours = aujourd'hui $-$ dateDebut
\end{itemize}
\end{formulebox}

\begin{exemplebox}[Exemple : ATG 1 — Bail 1 (mensuel)]
Loyer = 150\,000, charges = 7\,500, du 01/07/2023 au 31/08/2025 (792 jours)
$$\text{Attendu} = 157\,500 \times \frac{792}{30} = 4\,158\,000 \text{ FCFA}$$
\end{exemplebox}

\subsection{Bail ANNUEL}

\begin{formulebox}[Attendu — Bail annuel]
\begin{equation}
\text{Attendu} = (\text{loyer} + \text{charges}) \times 12 \times \text{nombre d'échéances passées}
\end{equation}
\vspace{-0.3cm}
\begin{itemize}[noitemsep]
  \item L'échéance tombe le \textbf{mois de début du bail}, chaque année
  \item On compte combien de ces mois anniversaires sont déjà passés
\end{itemize}
\end{formulebox}

\begin{exemplebox}[Exemple : ATG 1 — Bail 2 (annuel, début sept 2025)]
Au 03/06/2026 : échéance de septembre 2025 est passée $\Rightarrow$ 1 échéance.
$$\text{Attendu} = 165\,000 \times 12 \times 1 = 1\,980\,000 \text{ FCFA}$$
\end{exemplebox}

\subsection{Bail TRIMESTRIEL}

\begin{formulebox}[Attendu — Bail trimestriel]
\begin{equation}
\text{Attendu} = (\text{loyer} + \text{charges}) \times 3 \times \text{nombre d'échéances passées}
\end{equation}
\vspace{-0.3cm}
\begin{itemize}[noitemsep]
  \item L'échéance tombe tous les 3 mois à partir du mois de début
  \item Ex: début mai $\Rightarrow$ échéances en mai, août, novembre, février...
\end{itemize}
\end{formulebox}

\subsection{Bail SEMESTRIEL}

\begin{formulebox}[Attendu — Bail semestriel]
\begin{equation}
\text{Attendu} = (\text{loyer} + \text{charges}) \times 6 \times \text{nombre d'échéances passées}
\end{equation}
\end{formulebox}

% =================================================================
\section{Total Réglé}
% =================================================================

\begin{formulebox}[Total réglé]
\begin{equation}
\text{Réglé} = \sum_{\text{tous les baux du locataire sur l'appartement}} \sum_{\text{paiements}} \text{montant}
\end{equation}
\vspace{-0.3cm}
\begin{itemize}[noitemsep]
  \item Inclut les paiements de \textbf{tous} les baux (terminés + actif)
  \item N'inclut \textbf{pas} la caution (comptée séparément)
\end{itemize}
\end{formulebox}

% =================================================================
\section{Total Encaissé}
% =================================================================

\begin{formulebox}[Total encaissé]
\begin{equation}
\text{Encaissé} = \text{Réglé} + \sum_{\text{tous les baux}} \text{caution payée}
\end{equation}
\vspace{-0.3cm}
\begin{itemize}[noitemsep]
  \item La caution est comptée si \texttt{cautionPayee = true} sur le bail
  \item On somme les cautions de \textbf{tous} les baux (pas seulement l'actif)
\end{itemize}
\end{formulebox}

% =================================================================
\section{Différence (Avance / Impayé)}
% =================================================================

\begin{formulebox}[Différence]
\begin{equation}
\text{Différence} = \text{Réglé} - \text{Attendu}
\end{equation}
\vspace{-0.3cm}
\begin{itemize}[noitemsep]
  \item Positif $\Rightarrow$ le locataire est en \textbf{avance}
  \item Négatif $\Rightarrow$ le locataire a des \textbf{impayés}
  \item Zéro $\Rightarrow$ \textbf{à jour}
\end{itemize}
\end{formulebox}

% =================================================================
\section{Total Dû (Situation Globale)}
% =================================================================

Le total dû est calculé différemment de la différence simple. Il représente les \textbf{mois échus non réglés} du bail actif.

\subsection{Bail MENSUEL}

\begin{formulebox}[Total dû — mensuel]
Pour chaque mois passé depuis le début du bail :
\begin{equation}
\text{Si paiement du mois} < \text{loyer} + \text{charges} \Rightarrow \text{mois impayé}
\end{equation}
\begin{equation}
\text{Total dû} = \sum_{\text{mois impayés}} (\text{montant attendu} - \text{montant payé})
\end{equation}
\end{formulebox}

\subsection{Bail ANNUEL / TRIMESTRIEL / SEMESTRIEL}

\begin{formulebox}[Total dû — non-mensuel (proratisé)]
Pour une échéance impayée, on ne compte \textbf{pas} la totalité de l'échéance mais seulement les mois écoulés depuis :
\begin{equation}
\text{Total dû} = \text{totalMensuel} \times \min(\text{fréquence},\; \text{mois écoulés depuis l'échéance})
\end{equation}
\vspace{-0.3cm}
\begin{itemize}[noitemsep]
  \item La dette augmente \textbf{mois par mois}, pas d'un coup
  \item Plafonné à la fréquence (ex: max 12 pour un annuel)
\end{itemize}
\end{formulebox}

\begin{exemplebox}[Exemple : TMCO — Bail annuel, échéance jan 2026 impayée]
Au 03/06/2026 : 6 mois écoulés depuis janvier 2026.
$$\text{Total dû} = 157\,500 \times \min(12, 6) = 157\,500 \times 6 = 945\,000 \text{ FCFA}$$
\end{exemplebox}

\begin{exemplebox}[Exemple : BASS TECHNOLOGIES — Bail trimestriel, échéance mai 2026 impayée]
Au 03/06/2026 : 1 mois écoulé depuis mai 2026.
$$\text{Total dû} = 600\,000 \times \min(3, 1) = 600\,000 \text{ FCFA}$$
\end{exemplebox}

% =================================================================
\section{Détermination des Mois d'Échéance}
% =================================================================

\begin{formulebox}[Règle d'échéance]
Un mois $M$ est un mois d'échéance si :
\begin{equation}
(M_{\text{année}} - D_{\text{année}}) \times 12 + (M_{\text{mois}} - D_{\text{mois}}) \equiv 0 \pmod{\text{fréquence}}
\end{equation}
où $D$ = date de début du bail.

\vspace{0.2cm}
\begin{tabular}{ll}
\toprule
\textbf{Périodicité} & \textbf{Fréquence} \\
\midrule
Mensuel & 1 (tous les mois) \\
Trimestriel & 3 (tous les 3 mois) \\
Semestriel & 6 (tous les 6 mois) \\
Annuel & 12 (une fois par an) \\
\bottomrule
\end{tabular}
\end{formulebox}

\begin{exemplebox}[Exemples]
\begin{itemize}
  \item Bail ANNUEL débutant en \textbf{septembre} $\Rightarrow$ échéances en sept. chaque année
  \item Bail ANNUEL débutant en \textbf{janvier} $\Rightarrow$ échéances en janvier chaque année
  \item Bail TRIMESTRIEL débutant en \textbf{mai} $\Rightarrow$ échéances en mai, août, nov, fév
\end{itemize}
\end{exemplebox}

% =================================================================
\section{Renouvellement Automatique des Baux}
% =================================================================

\begin{formulebox}[Logique du CRON quotidien]
Chaque jour, pour chaque bail arrivant à expiration :
\begin{enumerate}[noitemsep]
  \item Vérifier \texttt{renouvellementAuto = true}
  \item Si oui : créer un nouveau bail avec les mêmes conditions
  \item Ancien bail $\rightarrow$ statut TERMINÉ
  \item Nouveau bail $\rightarrow$ statut ACTIF, mêmes montants, nouvelle période
\end{enumerate}
\textbf{Pas de condition sur les impayés} — le bail est renouvelé même si le locataire a des dettes.
\end{formulebox}

% =================================================================
\section{Reporting — Suivi des Paiements}
% =================================================================

\begin{formulebox}[Colonnes du tableau de reporting]
\begin{align}
\text{Jours d'habitation} &= \text{aujourd'hui} - \text{date du premier bail} \\
\text{Mois d'habitation} &= \frac{\text{jours d'habitation}}{30} \\
\text{Attendu} &= \text{formule multi-baux (§1)} \\
\text{Réglé} &= \text{total paiements tous baux (§2)} \\
\text{Différence} &= \text{Réglé} - \text{Attendu} \\
\text{Total encaissé} &= \text{Réglé} + \text{cautions payées (§3)}
\end{align}
\end{formulebox}

% =================================================================
\section{Finances — Revenus et Bilan}
% =================================================================

\begin{formulebox}[Calculs financiers]
\begin{align}
\text{Revenus du mois} &= \sum_{\text{paiements du mois}} \text{montant} \\
\text{Dépenses du mois} &= \sum_{\text{dépenses du mois}} \text{montant} \\
\text{Bilan net} &= \text{Revenus} - \text{Dépenses}
\end{align}
\end{formulebox}

% =================================================================
\section{Dashboard — Indicateurs}
% =================================================================

\begin{formulebox}[Formules des indicateurs]
\begin{align}
\text{Taux d'occupation} &= \frac{\text{appartements occupés}}{\text{total appartements}} \times 100 \\
\text{Total impayés} &= \sum_{\text{locataires}} \max(0,\; \text{Attendu} - \text{Réglé}) \\
\text{Revenus mois courant} &= \sum_{\text{paiements ce mois}} \text{montant}
\end{align}
\end{formulebox}

\end{document}
